\chapter[Aproximación metodológica]{Aproximación metodológica: del dato público al perfil semántico}
\label{ch:metodologia}

Este capítulo describe la aproximación propuesta en términos conceptuales, deliberadamente independientes de las tecnologías concretas empleadas (detallados en el Capítulo~\ref{ch:resultados}).
El propósito es doble: que el diseño sea evaluable por sí mismo y que resulte replicable por cualquier institución con otras herramientas equivalentes.

\section{Visión general}
\label{sec:met-vision}

El sistema se organiza en dos macrofases con ritmos distintos, ilustradas en la Figura~\ref{fig:flujo}:

\begin{itemize}
    \item \textbf{Preparación del corpus} (proceso periódico, fuera de línea): adquiere los datos públicos, los estructura, enriquece cada perfil y construye los índices semánticos.
    Se ejecuta de forma programada para incorporar la actividad investigadora nueva.
    \item \textbf{Atención de consultas} (en línea, tiempo real): recibe la idea libre del estudiante, la ``interpreta'', recupera los investigadores más afines y presenta una recomendación ordenada y justificada.
\end{itemize}

\begin{figure}[htbp]
    \centering
    \resizebox{\textwidth}{!}{
    \begin{tikzpicture}[
        font=\small,
        node distance=8mm and 7mm,
        etapa/.style={draw=upmblue!80!black, fill=upmblue!8, thick, rounded corners=3pt, align=center, minimum height=14mm, text width=31mm, inner sep=2mm},
        online/.style={draw=black!60, fill=black!4, thick, rounded corners=3pt, align=center, minimum height=14mm, text width=31mm, inner sep=2mm},
        flecha/.style={-{Stealth[length=2.8mm]}, semithick, draw=black!75}
    ]
    % Fase offline
    \node[etapa] (adq) {\textbf{Adquisición}\\[1mm] \footnotesize Censo y descarga de la actividad pública};
    \node[etapa, right=of adq] (est) {\textbf{Estructuración}\\[1mm] \footnotesize Grafo de investigadores y actividad};
    \node[etapa, right=of est] (enr) {\textbf{Enriquecimiento}\\[1mm] \footnotesize Biografías y dominios de especialización};
    \node[etapa, right=of enr] (ind) {\textbf{Indexación}\\[1mm] \footnotesize Representaciones dispersa y densa};

    \draw[flecha] (adq) -- (est);
    \draw[flecha] (est) -- (enr);
    \draw[flecha] (enr) -- (ind);

    % Fase online
    \node[online, below=22mm of adq] (idea) {\textbf{Idea del estudiante}\\[1mm] \footnotesize Texto libre, cualquier idioma};
    \node[online, below=22mm of est] (inter) {\textbf{Interpretación}\\[1mm] \footnotesize Extracción de términos o reformulación};
    \node[online, below=22mm of enr] (rec) {\textbf{Recuperación}\\[1mm] \footnotesize Similitud y fusión de rankings};
    \node[online, below=22mm of ind] (out) {\textbf{Recomendación}\\[1mm] \footnotesize Tutores ordenados, con evidencias};

    \draw[flecha] (idea) -- (inter);
    \draw[flecha] (inter) -- (rec);
    \draw[flecha] (rec) -- (out);
    \draw[flecha] (ind.south) -- node[right=-7mm, font=\footnotesize, align=left]{Espacio de\\representación\\compartido} (rec.north);

    % Agrupaciones
    \begin{scope}[on background layer]
    \node[draw=upmblue!50, dashed, rounded corners=4pt, fill=upmblue!3, inner sep=4mm,
        fit=(adq)(est)(enr)(ind),
        label={[upmblue!80!black, font=\footnotesize\bfseries]90:Preparación del corpus (proceso periódico)}] {};
    \node[draw=black!40, dashed, rounded corners=4pt, fill=black!2, inner sep=4mm,
        fit=(idea)(inter)(rec)(out),
        label={[black!70, font=\footnotesize\bfseries]270:Atención de consultas (tiempo real)}] {};
    \end{scope}
    \end{tikzpicture}
    }
    \caption{Flujo metodológico del sistema: preparación periódica del corpus (arriba) y atención de consultas en tiempo real (abajo).
    Ambas fases comparten el mismo espacio de representación semántica.}
    \label{fig:flujo}
\end{figure}

Tres principios de diseño atraviesan todo el flujo:

\begin{mydescription}
    \item[Etapas con artefactos persistentes.] Cada etapa consume el artefacto de la anterior y produce el suyo propio (páginas crudas, grafo, perfiles enriquecidos, índices).
    Esto permite re-ejecutar cualquier etapa de forma aislada e incremental (solo lo nuevo o lo inválido) sin repetir el proceso completo (\ref{req:adquisicion}).
    \item[Trazabilidad.] Toda recomendación es reconducible hacia atrás: del investigador recomendado a su perfil, del perfil a las evidencias agregadas y de estas a las páginas públicas originales.
    La fuente cruda se conserva inmutable precisamente para poder reprocesarla cuando el método mejora.
    \item[Neutralidad tecnológica.] Las etapas se definen por contratos de entrada/salida, no por herramientas. Cada una admite implementaciones alternativas (otro almacén, otro modelo de embedding) sin alterar el diseño.
\end{mydescription}

Conviene distinguir desde el principio dos artefactos que el flujo produce y que constituyen aportaciones diferenciadas, para que el lector no los confunda.
El \textbf{grafo de conocimiento} (etapa de estructuración) \emph{organiza} la actividad en entidades y relaciones (investigadores, publicaciones, proyectos, trabajos supervisados, asignaturas) y hace naturales las preguntas de conexión y vecindad.
La \textbf{representación semántica} (etapas de enriquecimiento e indexación) \emph{transforma} esa actividad en biografías, dominios de especialización y vectores localizables en un espacio de búsqueda.
En síntesis: el grafo dice \emph{cómo se relaciona} la evidencia y la representación semántica la hace \emph{buscable por afinidad}.

\section{Adquisición}
\label{sec:met-adquisicion}

La primera etapa convierte el portal público en una colección local y versionada y comprende dos pasos.
El primero es el \emph{censo}: construir un censo operativo de investigadores a partir de las vistas públicas disponibles en el portal, recorriendo el Portal Científico y recogiendo cada uno de los miembros.
El segundo es la \emph{descarga} de las dos vistas públicas de cada investigador, tanto el perfil (datos de filiación, posiciones, temas declarados), como su actividad (publicaciones, proyectos, trabajos supervisados).

La descarga se realiza de forma concurrente pero acotada, con reintentos ante fallos transitorios y de manera incremental: las páginas ya adquiridas no se vuelven a solicitar salvo refresco programado.
\duda{No sé si ser tran granular en la explicación}
El resultado se conserva tal cual se recibió, sin transformación alguna, como capa inmutable del sistema.


\section{Estructuración de la actividad en un grafo de conocimiento}
\label{sec:met-grafo}

La segunda etapa transforma el marcado semiestructurado en un modelo de datos explícito.
Se adopta un \emph{property graph} cuyos nodos representan las entidades del dominio (investigador, publicación, trabajo supervisado, posición en la universidad, asignatura impartida) y cuyas aristas capturan las relaciones entre ellas (autoría, supervisión, trabajo y enseñanza); el esquema completo de nodos y aristas se recoge en el Anexo~\ref{annex:modelo-datos}.

La elección del grafo frente a un modelo tabular responde a la naturaleza del dominio: las preguntas relevantes son de vecindad y conexión («¿qué investigadores comparten temas?», «¿quién ha supervisado trabajos próximos a este proyecto?») y el modelo de grafo las hace naturales.
Además, deja preparado el terreno para la línea futura de detección de colaboraciones entre investigadores, que es esencialmente un problema de análisis de grafos (Capítulo~\ref{ch:conclusiones}).

Durante esta etapa se aplica una limpieza \emph{conservadora} del texto: solo se transforma aquello que responde a patrones inequívocos de la fuente (por ejemplo, resúmenes con secciones etiquetadas concatenadas), prefiriéndose conservar una imperfección a introducir una corrupción silenciosa (requisito \ref{req:normalizacion}).

\section{Consolidación y enriquecimiento del perfil}
\label{sec:met-enriquecimiento}

La tercera etapa produce, para cada \gls{pdi}, un \emph{documento de perfil} autocontenido que agrega toda la evidencia pública disponible: su producción científica (publicaciones con sus resúmenes), proyectos recientes, trabajos supervisados, temas declarados, filiaciones y posiciones ocupadas y, cuando está disponible, su actividad docente (asignaturas impartidas).
Conviene recordar que un \gls{pdi} es Personal Docente \emph{e} Investigador: aunque la recomendación se apoya principalmente en la actividad investigadora y de supervisión, la estructura del perfil conserva también esa dimensión \emph{docente}, de modo que ambas vertientes queden representadas.
Este documento es la unidad de trabajo de todo lo que sigue.

Sobre él se aplica el enriquecimiento lingüístico mediante un modelo de lenguaje operado localmente, en coherencia con el requisito de despliegue en infraestructura propia (\ref{req:soberania}).
La producción científica se agrupa cronológicamente y cada año se sintetiza \emph{de forma autónoma}, sin arrastre de unos periodos a otros, en un par de artefactos complementarios:

\begin{itemize}
    \item Un \textbf{resumen anual} en prosa: una síntesis de aquello en lo que el investigador es experto a la luz de la producción de ese año.
    \item Una lista de \textbf{dominios de especialización} de ese año: áreas de conocimiento concisas extraídas de la misma evidencia.
    Dentro de un mismo año los dominios se acumulan con una regla de estabilidad (los ya establecidos se conservan y solo se añaden los genuinamente nuevos) que evita oscilaciones entre ejecuciones.
\end{itemize}

Trocear la trayectoria por años cumple un doble propósito: evita exceder la ventana de contexto del modelo en las carreras largas y produce una granularidad homogénea entre perfiles densos y dispersos.
Cada resumen anual se conserva como artefacto propio, lo que aporta trazabilidad sobre la evolución del investigador y deja la serie temporal preparada para que la etapa de indexación construya a partir de ella la biografía vigente (Sección~\ref{sec:met-indexacion}).

La generación es \emph{estructurada}: el modelo debe devolver objetos con campos tipados que se validan automáticamente y toda salida que no supere la validación se descarta y se reintenta.
El enriquecimiento queda persistido junto al resto del perfil, las etapas posteriores parten de estos resúmenes ya generados en lugar de reprocesar los textos originales.

\section{Representación e indexación semántica}
\label{sec:met-indexacion}

La cuarta etapa parte de la serie de resúmenes anuales y, para cada investigador, los condensa primero en una única \textbf{biografía vigente}.
La condensación es \emph{recurrente}: sobre una ventana temporal reciente se toma el año más antiguo como base y se pliega cada año posterior sobre el acumulado dando \emph{más peso al trabajo reciente}, de modo que la biografía refleje las líneas en las que el investigador trabaja hoy sin descartar la experiencia anterior que siga siendo pertinente.
Este esquema de resumen recurrente ponderado por recencia sigue la práctica revisada en el estado del arte para historiales ricos en texto (Sección~\ref{sec:sota-llm}).
La biografía vigente y sus dominios resultantes se persisten junto al perfil.

Sobre ellos se proyecta cada investigador a un espacio de representación buscable.
Atendiendo al análisis del Capítulo~\ref{ch:sota}, se mantienen deliberadamente \emph{dos} representaciones complementarias por investigador:

\begin{itemize}
    \item Una \textbf{representación léxica dispersa} de sus vocabularios temáticos (temas declarados y dominios extraídos), en la que cada componente es un término ponderado por su importancia contextual.
    Es eficiente e \emph{interpretable}, permite mostrar qué términos han producido cada coincidencia.
    Conservarla no responde solo a la transparencia: la señal léxica sigue siendo una línea base muy competitiva que, en ciertas tareas de recuperación, iguala o supera a los \emph{embeddings} densos (Sección~\ref{sec:sota-lexica}), de modo que mantenerla junto a la densa aporta robustez, no solo explicabilidad.
    \item Una \textbf{representación semántica densa} de su biografía vigente, que captura proximidad conceptual más allá de la superficie léxica y de la lengua en que se exprese la consulta.
\end{itemize}

Ambas se almacenan en índices de similitud que admiten búsqueda eficiente por vecindad y fusión nativa de resultados.
La consulta del estudiante se proyectará a estos mismos espacios para realizar la recomendación.

\section{Emparejamiento}
\label{sec:met-emparejamiento}

La fase en línea ofrece cuatro estrategias de emparejamiento, robustas al desajuste léxico y translingüe entre consulta y perfil (requisito \ref{req:emparejamiento}), que se corresponden con las familias revisadas en el estado del arte (Tabla~\ref{tab:posicionamiento}) y que el marco de evaluación comparará en igualdad de condiciones:

\begin{enumerate}
    \item \textbf{Léxica con interpretación.} Un modelo de lenguaje extrae de la idea del estudiante dos listas de términos en el registro de los perfiles (términos generales del área y términos específicos de método). Cada lista lanza una búsqueda dispersa y ambas se fusionan por RRF (Ecuación~\ref{eq:rrf}).
    \item \textbf{Semántica directa.} La idea, tal cual fue escrita, se proyecta al espacio denso y se recuperan las biografías más próximas por similitud de coseno.
    \item \textbf{Semántica enriquecida.} Antes de proyectar, un modelo de lenguaje reformula la idea: la traduce al registro técnico y la expande con terminología estrechamente relacionada, con la instrucción explícita de no introducir temas que el estudiante no haya pedido.
    \item \textbf{Híbrida.} Se lanzan en paralelo una señal léxica (términos extraídos contra los dominios de especialización) y una semántica (idea enriquecida contra las biografías) y ambas listas se fusionan por RRF.
\end{enumerate}

A estas cuatro estrategias del prototipo, el marco de evaluación añade una quinta de \emph{control} (léxica directa): lanza la búsqueda dispersa con el texto de la consulta tal cual, sin que el modelo de lenguaje extraiga ni reformule nada, lo que aísla la contribución del \gls{llm} y actúa como línea base clásica de coincidencia léxica (Sección~\ref{sec:eval-estrategias}).

En conjunto, el sistema emplea la \gls{rrf} (Ecuación~\ref{eq:rrf}) en dos puntos: para fusionar las dos listas de términos de la consulta en el modo léxico y para fusionar las señales léxica y semántica en el modo híbrido.
En todas ellas el resultado es una lista ordenada de investigadores con su puntuación normalizada y sus evidencias asociadas.

En todos los modos que invocan al modelo de lenguaje (la extracción de términos y la reformulación de la consulta) la generación se realiza con una \emph{temperatura de muestreo baja} ($t=0'1$).
Interpretar la idea es una tarea de transcripción fiel (traducir y acotar lo que el estudiante pide), no de redacción creativa: una temperatura baja favorece salidas estables y reproducibles ante la misma consulta y refuerza la prohibición de introducir temas que el estudiante no haya planteado.

\section{Presentación y transparencia}
\label{sec:met-presentacion}

La última etapa fija los requisitos de la capa de presentación.
El estudiante introduce su idea como texto libre, sin formularios ni vocabularios controlados y en su propio idioma.
La respuesta se presenta como una lista ordenada en la que cada recomendación va acompañada de sus \emph{evidencias}: los temas del investigador que coinciden con la consulta, su biografía investigadora y, cuando procede, la interpretación que el sistema hizo de la idea (términos extraídos o reformulación), de modo que el estudiante pueda auditar por qué se le recomienda a cada \gls{pdi}.

Este último punto es una decisión de diseño deliberada: el sistema se concibe como \emph{apoyo} a la decisión, nunca como sustituto.
Mostrar la interpretación de la consulta permite al estudiante corregirla, mostrar las evidencias le permite descartar recomendaciones espurias.
La evaluación de la calidad de estas recomendaciones, en términos objetivos, es el objeto del capítulo siguiente.
